% Source: contents/part-5-eng.tex (Figure: coupling restart mechanism / restart illustration)
\begin{tikzpicture}[
            xscale=2, 
            yscale=1.2, 
            >=Stealth,
            % 样式定义
            survivor/.style={circle, fill=blue!70!black, inner sep=1.2pt},
            killed-node/.style={circle, draw=blue!70!black, fill=white, inner sep=1.2pt, cross out, thick, minimum size=4pt},
            edge-live/.style={thick, blue!70!black},
            edge-killed/.style={thick, dashed, blue!30},
            cloud-tri/.style={fill=blue!10, opacity=0.5},
            ref-line/.style={thin, gray!40, dashed}
        ]

            % 1. 时间周期参考线 (0, l_N, 2l_N, 3l_N)
            \foreach \x/\label in {0/0, 2/\ell_N, 4/2\ell_N, 6/3\ell_N} {
                \draw[ref-line] (\x, -0.5) -- (\x, 4.5);
                \node[below] at (\x, -0.5) {$\label$};
            }

            % ================= 第一周期: 0 到 l_N =================
            % 起点
            \node[survivor] (n0) at (0, 0) {};
            \node[left] at (n0) {$\emptyset$};

            % 三角形阴影 (顶点在 n0)
            % 假设在 l_N 时刻，粒子分布在 [0.8, 2.2] 之间
            \fill[cloud-tri] (0,0) -- (2, 0.8) -- (2, 2.2) -- cycle;

            % 高频分支
            \draw[edge-live] (0,0) -- (0.5, 0.2) -- (1.0, 0.5) -- (1.5, 0.4) -- (2, 0.8);
            \draw[edge-killed] (0.5, 0.2) -- (1.2, 1.2) -- (2, 1.3);
            \draw[edge-killed] (1.0, 0.5) -- (1.6, 1.0) -- (2, 1.8);
            \draw[edge-killed] (0,0) -- (0.8, 0.9) -- (2, 2.2);

            % 瓶颈处粒子
            \node[survivor] (m1) at (2, 0.8) {}; 
            \node[killed-node] at (2, 1.3) {};
            \node[killed-node] at (2, 1.8) {};
            \node[killed-node] at (2, 2.2) {};

            % ================= 第二周期: l_N 到 2l_N =================
            % 起点是 m1 (2, 0.8)
            % 假设在 2l_N 时刻，粒子分布在 [1.8, 3.2] 之间
            \fill[cloud-tri] (2, 0.8) -- (4, 1.8) -- (4, 3.2) -- cycle;

            % 高频分支
            \draw[edge-live] (2, 0.8) -- (2.5, 1.1) -- (3.0, 1.4) -- (3.5, 1.3) -- (4, 1.8);
            \draw[edge-killed] (2.5, 1.1) -- (3.2, 2.0) -- (4, 2.3);
            \draw[edge-killed] (3.0, 1.4) -- (3.6, 2.1) -- (4, 2.8);
            \draw[edge-killed] (2, 0.8) -- (2.8, 2.2) -- (4, 3.2);

            % 瓶颈处粒子
            \node[survivor] (m2) at (4, 1.8) {};
            \node[killed-node] at (4, 2.3) {};
            \node[killed-node] at (4, 2.8) {};
            \node[killed-node] at (4, 3.2) {};

            % ================= 第三周期: 2l_N 到 3l_N =================
            % 起点是 m2 (4, 1.8)
            % 假设在 3l_N 时刻，粒子分布在 [2.7, 4.1] 之间
            \fill[cloud-tri] (4, 1.8) -- (6, 2.7) -- (6, 4.1) -- cycle;

            % 高频分支
            \draw[edge-live] (4, 1.8) -- (4.5, 2.0) -- (5.0, 2.4) -- (5.5, 2.3) -- (6, 2.7);
            \draw[edge-killed] (4.5, 2.0) -- (5.2, 2.9) -- (6, 3.2);
            \draw[edge-killed] (5.0, 2.4) -- (5.6, 3.1) -- (6, 3.7);
            \draw[edge-killed] (4, 1.8) -- (4.8, 3.2) -- (6, 4.1);

            % 瓶颈处粒子
            \node[survivor] (m3) at (6, 2.7) {};
            \node[killed-node] at (6, 3.2) {};
            \node[killed-node] at (6, 3.7) {};
            \node[killed-node] at (6, 4.1) {};

            % 瓶颈标注
            \node[above right] at (m1) {\scriptsize keep leftmost};
            \node[above right] at (m2) {\scriptsize keep leftmost};
            \node[above right] at (m3) {\scriptsize keep leftmost};
\end{tikzpicture}

